% !TeX root = ../../Book2.tex
% 纲要标题：### 1.3 线性算子
% 纲要位置：计划纲要.md，第 843 行。

\section{线性算子}
\label{b2:sec:0103}

给定一个向量及一个线性算子，可以反复施加算子，得到 \(v,Tv,T^2v,\ldots\)。这些向量之间的关系记录了算子的一部分结构。本节把关系写成多项式，说明一个算子如何使整个多项式环作用在空间上。这将是下一章从向量空间转向模的一个主要例子。

% 纲要内容（仅作写作计划，尚非教材正文）：
%
% * 特征子空间、广义特征子空间及互素核分解
% * 最小多项式与循环向量
% * 将线性算子看成多项式环的作用；限制与商算子的两层零化关系
%

\subsection{特征子空间与广义特征子空间}
\label{b2:subsec:0103:01}

\begin{definition}\label{b2:def:linear-operator}
设 \(V\) 是域 \(k\) 上的向量空间。给定线性映射 \(T:V\to V\)，称 \(T\) 为 \(V\) 上的\term[xianxing-suanzi]{线性算子}[linear operator]，也称线性自同态。记
\(\operatorname{End}_k(V)=\operatorname{Hom}_k(V,V)\)。其中加法逐点定义，乘法采用复合，约定 \(ST=S\circ T\)，即先施加右侧的 \(T\)。恒等算子记为 \(I_V\)，并令 \(T^0=I_V\)。
\end{definition}

\symidx{\(\operatorname{End}_k(V)\)：向量空间 \(V\) 的线性自同态环}

复合满足结合律，又对逐点加法满足左右分配律，所以 \(\operatorname{End}_k(V)\) 是含幺环。它通常不交换；标量算子 \(aI_V\) 则与每个线性算子交换，因为 \(T(av)=a\cdot T(v)\)。以下除单独讨论外，\(V\) 有限维且非零，记 \(n=\dim_kV\)。

\begin{definition}\label{b2:def:eigen-generalized}
设 \(T\) 是域 \(k\) 上的向量空间 \(V\) 的线性算子，\(\lambda\in k\)。若 \(0\ne v\in V\) 且 \(Tv=\lambda v\)，就称 \(v\) 为 \(T\) 对应于 \(\lambda\) 的\term[tezheng-xiangliang]{特征向量}[eigenvector]，并称 \(\lambda\) 为 \(T\) 的\term[tezheng-zhi]{特征值}[eigenvalue]。对应于 \(\lambda\) 的\term[tezheng-zi-kongjian]{特征子空间}[eigenspace]和\term[guangyi-tezheng-zi-kongjian]{广义特征子空间}[generalized eigenspace]为
\[
E_\lambda(T)=\ker(T-\lambda I_V),\qquad
G_\lambda(T)=\bigcup_{r\ge1}\ker\bigl((T-\lambda I_V)^r\bigr).
\]
\end{definition}
\symidx{\(E_\lambda(T),G_\lambda(T)\)：特征子空间与广义特征子空间}

这些核依次包含，故其并是子空间。每个核都被 \(T\) 保持：\(T\) 与 \((T-\lambda I_V)^r\) 交换。一般地，设 \(T\) 是 \(k\) 向量空间 \(V\) 上的算子。若子空间 \(U\le V\) 满足 \(T(U)\subseteq U\)，就称 \(U\) 为 \(T\) 的\term[bubian-zi-kongjian]{不变子空间}[invariant subspace]。

\begin{proposition}\label{b2:prop:kernel-power-stability}
设 \(V\) 是域 \(k\) 上的 \(n\) 维向量空间，\(N\in\operatorname{End}_k(V)\)，则核链 \(\ker(N^r)\) 至迟在 \(r=n\) 时稳定，即 \(\ker(N^n)=\ker(N^{n+1})=\cdots\)。
\end{proposition}
\begin{proof}
令 \(K_r=\ker(N^r)\)，其中 \(K_0=0\)。若 \(K_r=K_{r+1}\)，对 \(v\in K_{r+2}\) 有 \(Nv\in K_{r+1}=K_r\)，所以 \(v\in K_{r+1}\)。因此一旦相邻两项相等，以后各项都相等。每一次严格包含至少使维数增加一；从维数零开始至多增加 \(n\) 次，故 \(K_n=K_{n+1}=\cdots\)。
\end{proof}

将命题中的 \(N\) 取为 \(T-\lambda I_V\)，就得到
\[
\begin{gathered}
\ker(T-\lambda I_V)\subseteq\ker\bigl((T-\lambda I_V)^2\bigr)\subseteq\cdots,\\
G_\lambda(T)=\ker\bigl((T-\lambda I_V)^n\bigr).
\end{gathered}
\]
因此在有限维空间中，不必取无穷多个核的并，只需计算核链稳定后的一项。

\ProseParagraphBreak
例如在 \(k^2\) 的基 \(e_1,e_2\) 下取
\[
Te_1=\lambda e_1,\qquad Te_2=e_1+\lambda e_2.
\]
此时 \(N=T-\lambda I_V\) 满足 \(Ne_2=e_1,Ne_1=0\)。所以 \(E_\lambda=ke_1\)，而 \(G_\lambda=V\)。对应于 \(\lambda\) 的广义特征向量，记录反复施加 \(T-\lambda I_V\) 后何时变为零；普通特征向量在第一次施加该算子时便变为零。

\ProseParagraphBreak
不同特征值的特征向量线性无关。确切地说，若存在非平凡关系，取其中非零项数最少的一条，写成
\[
\sum_{j=1}^ra_jv_j=0,
\]
其中各 \(\lambda_j\) 不同且 \(a_j\ne0\)，则 \(r\ge2\)。施加 \(T-\lambda_rI_V\)，得到少一项的关系，其余系数 \(a_j(\lambda_j-\lambda_r)\) 仍非零，矛盾。不过特征向量未必生成全空间，上面的二维例子已经说明这一点；基域中还可能根本没有特征值。

\subsection{最小多项式与循环向量}
\label{b2:subsec:0103:02}

对于 \(p(t)=\sum_{i=0}^da_it^i\in k[t]\)，定义
\[
p(T)=\sum_{i=0}^da_iT^i.
\]
标量算子与 \(T\) 交换，因而直接展开有限和可得
\[
(p+q)(T)=p(T)+q(T),\qquad (pq)(T)=p(T)q(T),\qquad 1(T)=I_V.
\]
这些公式说明求值映射
\[
\Phi_T:k[t]\longrightarrow\operatorname{End}_k(V),\qquad p\longmapsto p(T)
\]
保持加法、乘法和单位。记它的像为
\[
k[T]\coloneqq\operatorname{im}\Phi_T=\{\,p(T)\mid p\in k[t]\,\}.
\]
\symidx{\(k[T]\)：由线性算子生成的多项式算子代数}
这里 \(k[t]\) 中的元素是形式多项式，\(k[T]\) 中的元素则是实际作用于 \(V\) 的算子；不同多项式可能给出同一个算子，例如 \(T^2=0\) 时 \(t^2\) 与零多项式的像相同。虽然 \(\operatorname{End}_k(V)\) 通常不交换，\(k[T]\) 中的算子却彼此交换。

\begin{definition}\label{b2:def:minimal-polynomial}
设 \(V\) 是域 \(k\) 上的向量空间，\(T\in\operatorname{End}_k(V)\)，\(p\in k[t]\)。若 \(p(T)=0\)，就称 \(p\) \term[linghua]{零化}[annihilate]算子 \(T\)。如果存在非零零化多项式，就把其中次数最小的首一多项式称为 \(T\) 的\term[zuixiao-duoxiangshi]{最小多项式}[minimal polynomial]，记作 \(\mu_T(t)\)。
\end{definition}
\symidx{\(\mu_T\)：线性算子的最小多项式}

零向量空间上的恒等算子也是零算子，因此这时 \(\mu_T=1\)。

\begin{proposition}\label{b2:prop:minimal-polynomial-divisibility}
设 \(V\) 是域 \(k\) 上的有限维向量空间，\(T\in\operatorname{End}_k(V)\)，则 \(T\) 有唯一的最小多项式 \(\mu_T\)，且对任意 \(p\in k[t]\) 有
\[
p(T)=0\quad\Longleftrightarrow\quad \mu_T\mid p.
\]
\end{proposition}
\begin{proof}
若 \(V=0\)，任意 \(p(T)\) 都是零算子，而 \(\mu_T=1\)，结论立即成立。以下设 \(V\ne0\)，令 \(n=\dim_kV\)。在一组基上指定像，便以 \(n^2\) 个标量确定一个自同态，所以 \(\operatorname{End}_k(V)\) 维数为 \(n^2\)。因此 \(I_V,T,\ldots,T^{n^2}\) 线性相关，存在非零零化多项式。选次数最小者，再除以其首项系数，得到首一的 \(\mu_T\)。它的次数至少为一，因为非零常数算子不是零算子。

对任意 \(p\) 作带余除法 \(p=qm_T+r\)，其中 \(r=0\) 或 \(\deg r<\deg \mu_T\)。若 \(p(T)=0\)，则 \(r(T)=0\)；最小性迫使 \(r=0\)。反向蕴含由乘法公式立即得到。另一同次数的首一最小零化多项式必须为 \(\mu_T\) 的常数倍，首一条件又使这个倍数为一。
\end{proof}

上述证明中的 \(n^2\) 只给出存在非零零化多项式的粗界。
\cref{b2:prop:operator-characteristic-minimal}将在第四章证明 \(\mu_T\mid\chi_T\)，从而得到更强的 \(\deg \mu_T\le n\)。这里也已确定求值映射的核：\(\ker\Phi_T=(\mu_T)\)，即 \(p(T)=q(T)\) 当且仅当 \(\mu_T\mid(p-q)\)。

\begin{corollary}\label{b2:cor:eigenvalue-minimal-root}
设 \(V\ne0\) 是域 \(k\) 上的有限维向量空间，\(T\in\operatorname{End}_k(V)\)，\(\lambda\in k\)。则 \(\lambda\) 是 \(T\) 的特征值，当且仅当 \(\mu_T(\lambda)=0\)。
\end{corollary}
\begin{proof}
若 \(0\ne v\in V\) 满足 \(Tv=\lambda v\)，则 \(0=\mu_T(T)v=\mu_T(\lambda)v\)，故 \(\mu_T(\lambda)=0\)。反过来，若 \(\mu_T=(t-\lambda)q\) 而 \(T-\lambda I_V\) 单射，由 \((T-\lambda I_V)q(T)=0\) 可得 \(q(T)=0\)，与 \(\mu_T\) 的最小性矛盾。因此 \(T-\lambda I_V\) 有非零核。
\end{proof}

\begin{definition}\label{b2:def:vector-minimal-polynomial}
设 \(V\) 是域 \(k\) 上的向量空间，\(T\in\operatorname{End}_k(V)\)，\(v\in V\)。若存在非零多项式 \(p\in k[t]\) 使 \(p(T)v=0\)，则在这样的多项式中取次数最小的首一者，称为 \(v\) 相对于 \(T\) 的最小多项式，记作 \(\mu_{T,v}\)。特别地，\(\mu_{T,0}=1\)。
\end{definition}
\symidx{\(\mu_{T,v}\)：向量相对于算子的最小多项式}

带余除法表明 \(\mu_{T,v}\) 唯一，且 \(p(T)v=0\) 当且仅当 \(\mu_{T,v}\mid p\)。若 \(V\) 有限维，对非零 \(v\) 可利用 \(v,Tv,\ldots,T^{\dim_kV}v\) 的线性相关性得到非零零化多项式，因此每个向量都有 \(\mu_{T,v}\)。只要 \(\mu_T\) 存在，\(\mu_T(T)v=0\) 还给出 \(\mu_{T,v}\mid \mu_T\)。

\begin{definition}\label{b2:def:cyclic-vector}
设 \(V\) 是域 \(k\) 上的向量空间，\(T\in\operatorname{End}_k(V)\)，\(v\in V\)。由 \(v\) 产生的\term[xunhuan-zi-kongjian]{循环子空间}[cyclic subspace]是
\[
k[T]v=\{p(T)v\mid p\in k[t]\}.
\]
它是所有 \(T^jv\ (j\ge0)\) 的线性生成空间，也是 \(k[T]\) 中的算子作用于 \(v\) 所得到的全部向量组成的空间。若 \(k[T]v=V\)，就称 \(v\) 为 \(T\) 的\term[xunhuan-xiangliang]{循环向量}[cyclic vector]。
\end{definition}
\symidx{\(k[T]v\)：由向量生成的循环子空间}

映射 \(k[t]\to k[T]v\)、\(p\mapsto p(T)v\) 的核记录了 \(v\) 满足的全部多项式关系。有非零关系时，这个核是 \((\mu_{T,v})\)；没有非零关系时，核为零。与 \(\Phi_T\) 的核不同，这里只要求零化 \(v\)，不要求零化整个空间。

\begin{proposition}\label{b2:prop:cyclic-basis}
设 \(V\ne0\) 是域 \(k\) 上的有限维向量空间，\(T\in\operatorname{End}_k(V)\)，\(0\ne v\in V\)。令 \(\mu_{T,v}\) 为零化 \(v\) 的首一最小多项式，\(d=\deg \mu_{T,v}\)。此时 \(v,Tv,\ldots,T^{d-1}v\) 是 \(k[T]v\) 的基。若 \(v\) 是循环向量，则 \(\mu_T=\mu_{T,v}\)，从而 \(\deg \mu_T=\dim_k V\)。
\end{proposition}
\begin{proof}
这些向量若相关，就有一个次数小于 \(d\) 的非零多项式零化 \(v\)，矛盾。对任意 \(p\)，用 \(\mu_{T,v}\) 除去高次项，得到 \(p(T)v=r(T)v\)、\(\deg r<d\)，所以它们生成循环子空间。

若这个空间为 \(V\)，则每个向量写成 \(p(T)v\)。算子多项式彼此交换给出
\[
\mu_{T,v}(T)p(T)v=p(T)\mu_{T,v}(T)v=0,
\]
故 \(\mu_{T,v}\) 零化整个算子，上一命题给出 \(\mu_T\mid \mu_{T,v}\)。另一方面，\(\mu_T(T)v=0\) 给出 \(\mu_{T,v}\mid \mu_T\)。两者均首一，故相等。
\end{proof}

循环向量不总存在。\(n>1\) 时，标量算子 \(T=\lambda I_V\) 的每个非零循环子空间都只有一维。相反，前面的二维例子以 \(e_2\) 为循环向量，因为 \(e_2,Te_2\) 生成 \(e_1,e_2\)；其最小多项式为 \((t-\lambda)^2\)。

\begin{proposition}[互素多项式给出的分解]\label{b2:prop:coprime-operator-decomposition}
设 \(V\) 是域 \(k\) 上的向量空间，\(T\in\operatorname{End}_k(V)\)。若 \(a,b\in k[t]\) 互素且 \(a(T)b(T)=0\)，则
\[
V=\ker a(T)\oplus\ker b(T).
\]
\end{proposition}
\begin{proof}
由第一卷第11.2节的 Euclid 算法命题\footnote{欧几里得（希腊文 \OriginalName{Εὐκλείδης}，约公元前325-约公元前265），古希腊数学家，著有《几何原本》，其中系统讨论了辗转相除法。}给出的 Bézout 等式\footnote{贝祖（Étienne Bézout，1730-1783），法国数学家，研究代数方程的消元与结式，并著有数学教材。}，互素多项式 \(a,b\) 的最大公因子为单位，故可取 \(r,s\in k[t]\)，使 \(ra+sb=1\)。于是
\[
v=r(T)a(T)v+s(T)b(T)v.
\]
第一项被 \(b(T)\) 零化，第二项被 \(a(T)\) 零化，故两核之和为全空间。若 \(v\) 同时在两核中，同一等式给出 \(v=0\)，所以和是直和。核由算子多项式定义，也都是 \(T\) 不变子空间。
\end{proof}

例如 \((T-\lambda I_V)(T-\mu I_V)=0\) 且 \(\lambda\ne\mu\) 时，每个向量唯一分成两个特征子空间中的分量。最小多项式中出现重复的一次因子时，同一互素分解把普通特征子空间换成广义特征子空间。

\begin{corollary}\label{b2:cor:generalized-eigenspace-decomposition}
设 \(V\ne0\) 是域 \(k\) 上的有限维向量空间，\(T\in\operatorname{End}_k(V)\)。若
\[
\mu_T(t)=\prod_{i=1}^s(t-\lambda_i)^{a_i},
\qquad \lambda_i\in k\;\text{两两不同},\quad a_i\ge1,
\]
则
\[
V=\bigoplus_{i=1}^sG_{\lambda_i}(T),\qquad
G_{\lambda_i}(T)=\ker\bigl((T-\lambda_iI_V)^{a_i}\bigr).
\]
各个直和分量都是 \(T\) 不变子空间。
\end{corollary}
\begin{proof}
令 \(q_i=(t-\lambda_i)^{a_i}\)。这些多项式两两互素，且 \(\prod_iq_i(T)=\mu_T(T)=0\)。先证明一般的直和公式：只要 \(q_1,\ldots,q_s\) 两两互素且乘积零化算子，就有 \(V=\bigoplus_i\ker q_i(T)\)。对因子个数 \(s\) 归纳：\(s=1\) 时 \(q_1(T)=0\)；\(s>1\) 时先将 \(q_1\) 与 \(q_2\cdots q_s\) 代入\cref{b2:prop:coprime-operator-decomposition}，得到
\[
V=\ker q_1(T)\oplus\ker\bigl(q_2(T)\cdots q_s(T)\bigr).
\]
第二个分量被 \(T\) 保持，在其上应用归纳假设便得到其余各核的直和。这里每个 \(\ker q_i(T)\ (i\ge2)\) 已包含在第二个分量中，所以所得分解为 \(V=\bigoplus_i\ker q_i(T)\)。

记 \(V_i=\ker q_i(T)\)，则定义直接给出 \(V_i\subseteq G_{\lambda_i}(T)\)。反过来，取 \(v\in G_{\lambda_i}(T)\)，有某个 \(r\ge1\) 使 \((T-\lambda_iI_V)^rv=0\)。沿上述直和写 \(v=\sum_jv_j\)，其中 \(v_j\in V_j\)。各分量被 \(T\) 保持，故直和的唯一性给出 \((T-\lambda_iI_V)^rv_j=0\)。当 \(j\ne i\) 时，\(q_j\) 与 \((t-\lambda_i)^r\) 互素，Bézout 等式表明一个向量若同时被这两个多项式零化，就只能为零。因此 \(v_j=0\ (j\ne i)\)，于是 \(v\in V_i\)。
\end{proof}

这个分解先将不同特征值的部分分开；在每个 \(G_{\lambda_i}(T)\) 内，\(T-\lambda_iI_V\) 是幂零算子。幂零部分还能怎样分成循环子空间，是 Jordan 块需要进一步描述的问题。

\subsection{线性算子的多项式环作用}
\label{b2:subsec:0103:03}

现在把 \(T\) 连同全部 \(p(T)\) 一起考虑。在 \(V\) 原有的加法上规定
\[
p(t)\cdot v=p(T)v.
\]
算子多项式的加法与乘法公式给出
\[
(p+q)\cdot v=p\cdot v+q\cdot v,\qquad
p\cdot(u+v)=p\cdot u+p\cdot v,\qquad
(pq)\cdot v=p\cdot(q\cdot v),\qquad 1\cdot v=v.
\]
常数多项式的作用与原来的 \(k\) 标量乘法一致。下一章会把这种结构称为 \(k[t]\) 上的模。对于任意维空间，以上定义仍成立；只是算子未必有非零零化多项式。例如在 \(k[t]\) 上令 \(T\) 为乘以 \(t\)，若 \(p(T)=0\)，把它作用在 \(1\) 上即得 \(p(t)=0\)。

\begin{proposition}\label{b2:prop:polynomial-action-subspace}
设 \(V,W\) 是域 \(k\) 上的向量空间，\(U\le V\)，\(T\in\operatorname{End}_k(V)\)、\(S\in\operatorname{End}_k(W)\)，则 \(U\) 被所有 \(p(T)\)（\(p\in k[t]\)）保持，当且仅当它被 \(T\) 保持。对任意 \(k\) 线性映射 \(F:V\to W\)，有
\[
FT=SF
\quad\Longleftrightarrow\quad
F(p(T)v)=p(S)F(v)\quad\text{对所有}\;p\in k[t],v\in V.
\]
\end{proposition}
\begin{proof}
若 \(T(U)\subseteq U\)，反复施加 \(T\) 后仍在 \(U\)，再取有限线性组合就得到 \(p(T)(U)\subseteq U\)；反向取 \(p=t\)。若 \(FT=SF\)，归纳给出 \(FT^j=S^jF\)，包括 \(j=0\)，再由线性性对各项相加即可。反向同样取 \(p=t\)。
\end{proof}

因此，\(T\) 不变子空间恰好是对全部 \(k[t]\) 作用封闭的子空间，下一章称它为子模；关系 \(FT=SF\) 恰好表示 \(F\) 保持 \(k[t]\) 作用，因而是 \(k[t]\) 线性映射。循环子空间 \(k[T]v\) 由一个向量经这种作用生成，称为循环模，而 \(p(T)v=0\) 是这个生成元的关系。这就把不变子空间、算子之间的映射和循环部分的关系放在同一种标量作用下研究。

\ProseParagraphBreak
这种作用也能限制到不变子空间，并下降到商空间。若一个多项式已经零化商算子，它在原空间上的像就落入所取的子空间；再用零化限制算子的多项式，便能零化整个空间。

\begin{proposition}\label{b2:prop:restriction-quotient-polynomials}
设 \(V\) 是域 \(k\) 上的有限维向量空间，\(T\in\operatorname{End}_k(V)\)，\(U\le V\) 满足 \(T(U)\subseteq U\)。则限制 \(T_U=T|_U\) 是 \(U\) 上的算子，且
\[
\bar T:V/U\longrightarrow V/U,\qquad v+U\longmapsto Tv+U
\]
定义了商空间上的算子。商映射 \(q:V\to V/U\) 满足 \(qT=\bar Tq\)，并且
\[
\operatorname{lcm}(\mu_{T_U},\mu_{\bar T})\mid \mu_T
\quad\text{且}\quad \mu_T\mid \mu_{T_U}\mu_{\bar T},
\]
其中最小公倍式取首一者，零空间上的最小多项式约定为 \(1\)。
\end{proposition}
\begin{proof}
若 \(v+U=w+U\)，则 \(v-w\in U\)，不变性给出 \(Tv-Tw\in U\)，所以 \(\bar T\) 与代表元无关。商空间的运算又给出其线性性，且 \(\bar Tq(v)=Tv+U=qT(v)\)。由\cref{b2:prop:polynomial-action-subspace}，对每个 \(p\in k[t]\) 都有
\[
q\,p(T)=p(\bar T)\,q.
\]
若 \(p(T)=0\)，则限制到 \(U\) 后有 \(p(T_U)=0\)，又由上式及 \(q\) 的满射性得到 \(p(\bar T)=0\)。取 \(p=\mu_T\)，\cref{b2:prop:minimal-polynomial-divisibility}给出 \(\mu_{T_U}\mid \mu_T\) 与 \(\mu_{\bar T}\mid \mu_T\)，从而它们的首一最小公倍式也整除 \(\mu_T\)。

令 \(a=\mu_{T_U}\)、\(b=\mu_{\bar T}\)。由于 \(b(\bar T)=0\)，有 \(q\,b(T)=0\)，即 \(b(T)V\subseteq U\)。而 \(a(T)\) 在 \(U\) 上等于 \(a(T_U)=0\)，所以 \(a(T)b(T)=0\)，即 \((ab)(T)=0\)。再用\cref{b2:prop:minimal-polynomial-divisibility}便得 \(\mu_T\mid ab\)。
\end{proof}

乘积一般不能换成最小公倍式。例如在 \(k^2\) 上令 \(Te_1=0\)、\(Te_2=e_1\)，并取 \(U=ke_1\)。限制与商算子都为零，其最小多项式都是 \(t\)；但 \(T\ne0\)、\(T^2=0\)，故 \(\mu_T=t^2\)。子空间和商空间分别只需作用一次就变为零，原空间却可能先从商空间落入子空间，再作用一次才变为零。
